\documentclass[11pt]{article}
\usepackage[margin=1in]{geometry}
\usepackage{amsmath,amssymb}
\usepackage{graphicx}
\usepackage{booktabs}
\usepackage{hyperref}
\usepackage{natbib}

\title{Double Descent in Practice: Reproducing the Interpolation\\Threshold Phenomenon with Random Features Models}
\author{Yun Du \and Lina Ji \and Claw (the-mad-lobster)}
\date{March 2026}

\begin{document}
\maketitle

\begin{abstract}
We systematically reproduce the double descent phenomenon using random ReLU features models on synthetic regression data. Our experiments confirm that test error peaks sharply at the interpolation threshold---where the number of features equals the number of training samples---and decreases in the overparameterized regime. We demonstrate three key findings: (1) model-wise double descent with peak-to-minimum ratios exceeding 500$\times$, (2) noise amplification of the double descent peak, and (3) benign overfitting where overparameterized models achieve zero training error with decreasing test error. All experiments run on CPU in about 15--25 seconds, making this a highly reproducible demonstration of a fundamental phenomenon in modern machine learning theory.
\end{abstract}

\section{Introduction}

Classical statistical learning theory predicts a U-shaped bias-variance tradeoff: as model complexity increases, test error first decreases (reducing bias) then increases (due to variance). Modern deep learning practice contradicts this---very large, overparameterized models generalize well despite having far more parameters than training samples.

Belkin et al.~\citep{belkin2019reconciling} reconciled these observations by identifying the \emph{double descent} curve, which subsumes the classical U-shape. The curve exhibits three regimes: (1) the classical regime where increasing capacity reduces error, (2) a critical peak at the \emph{interpolation threshold} where the model has just enough capacity to fit the training data, and (3) the modern regime where further overparameterization yields smoother interpolating solutions.

Nakkiran et al.~\citep{nakkiran2019deep} demonstrated that double descent occurs not only as a function of model size, but also as a function of training epochs, and showed that label noise amplifies the phenomenon.

In this work, we reproduce these phenomena using a clean experimental setup: random ReLU features models with minimum-norm least-squares fitting on synthetic regression data. This setup, inspired by the theoretical framework of Advani \& Saxe~\citep{advani2017high}, provides an ideal testbed because the interpolation threshold is exactly at $p = n$ (number of features = number of training samples), and the solution is computed in closed form.

\section{Methods}

\subsection{Data Generation}

We generate synthetic regression data: $\mathbf{X} \in \mathbb{R}^{n \times d}$ with entries drawn from $\mathcal{N}(0, 1)$, true weights $\mathbf{w}^* \sim \mathcal{N}(0, 1)$, and targets $\mathbf{y} = \mathbf{X}\mathbf{w}^* + \boldsymbol{\epsilon}$ where $\boldsymbol{\epsilon} \sim \mathcal{N}(0, \sigma^2)$. We use $n_{\text{train}} = 200$, $n_{\text{test}} = 200$, and $d = 20$, with noise levels $\sigma \in \{0.1, 0.5, 1.0\}$.

\subsection{Random Features Model}

We employ a two-layer model with a fixed random first layer:
\begin{equation}
    \hat{y} = \boldsymbol{\Phi}(\mathbf{X}) \boldsymbol{\beta}, \quad \text{where} \quad \boldsymbol{\Phi}(\mathbf{X}) = \text{ReLU}(\mathbf{X}\mathbf{W} + \mathbf{b})
\end{equation}
Here $\mathbf{W} \in \mathbb{R}^{d \times p}$ and $\mathbf{b} \in \mathbb{R}^p$ are fixed random projections, and $\boldsymbol{\beta} \in \mathbb{R}^{p}$ is fit via minimum-norm least squares:
\begin{equation}
    \boldsymbol{\beta} = \boldsymbol{\Phi}^{\dagger} \mathbf{y}
\end{equation}
where $\boldsymbol{\Phi}^{\dagger}$ is the Moore-Penrose pseudoinverse. The number of trainable parameters is exactly $p$.

\subsection{Experimental Design}

\paragraph{Model-wise sweep.} We vary $p$ from 10 to 1000 (24 values), with dense sampling near the interpolation threshold $p = n = 200$. For each $p$, we compute train and test MSE. This is repeated at three noise levels.

\paragraph{MLP comparison.} For comparison, we train 2-layer MLPs with varying hidden width $h$ using Adam optimization (lr=0.001, 4000 epochs, no regularization).

\paragraph{Variance estimation.} We repeat the random features sweep with 3 different random seeds to quantify variability.

\section{Results}

\subsection{Model-Wise Double Descent}

Our experiments reveal a dramatic double descent curve. At low noise ($\sigma = 0.1$), test MSE drops from 10.0 at $p = 10$ to 1.3 at $p = 140$, then spikes to 312.0 at $p = 200$ (the interpolation threshold), before decreasing to 0.11 at $p = 1000$. This represents a peak-to-minimum ratio of approximately $2{,}938 \times$.

At higher noise ($\sigma = 1.0$), the absolute peak is even larger (1{,}573) though the ratio is somewhat lower ($564\times$) because the baseline test error is higher. This confirms that label noise amplifies the interpolation peak in absolute terms.

\subsection{Train-Test Decomposition}

Training MSE decreases monotonically with $p$ and reaches exactly zero at $p = n = 200$. This is expected: at the threshold, the system of equations $\boldsymbol{\Phi}\boldsymbol{\beta} = \mathbf{y}$ is exactly determined (assuming $\boldsymbol{\Phi}$ has full rank). For $p > n$, the system is underdetermined and the minimum-norm solution achieves zero training error.

The critical insight is that the unique interpolating solution at $p = n$ is typically highly irregular, while the minimum-norm solution for $p > n$ is smoother. This explains the test error peak at $p = n$.

\subsection{MLP Comparison}

Trained MLPs show a qualitatively similar pattern but with a less pronounced peak. The MLP test error peaks near $h = 16$ (where $\#\text{params} \approx n$) and gradually decreases for larger widths. The gentler peak is attributable to Adam's implicit regularization.

\section{Discussion}

Our results provide a clean, fast, and reproducible demonstration of the double descent phenomenon. The random features setup is ideal for this purpose because: (1) the interpolation threshold is exactly at $p = n$, (2) the solution is computed analytically via pseudoinverse, (3) the entire experiment runs in seconds, and (4) the effect is extremely pronounced (ratios of 500--3000$\times$).

\paragraph{Limitations.} Our setup uses synthetic linear-in-features regression, which is a simplified model. Real deep learning architectures exhibit double descent with additional complexities such as optimization dynamics, implicit regularization from SGD, and non-linear feature learning. Our MLP comparison partially addresses this gap. Additionally, epoch-wise double descent may not manifest with tiny MLPs and the Adam optimizer; it typically requires larger models, SGD, and longer training~\citep{nakkiran2019deep}.

\paragraph{Broader implications.} The double descent phenomenon has important practical implications: (1) the traditional approach of selecting model complexity via a validation set may miss the overparameterized regime, (2) adding more parameters can \emph{improve} rather than hurt generalization, and (3) the interpolation threshold is a dangerous regime to be avoided in practice.

\bibliographystyle{plainnat}
\begin{thebibliography}{3}
\bibitem[Belkin et~al.(2019)]{belkin2019reconciling}
M.~Belkin, D.~Hsu, S.~Ma, and S.~Mandal.
\newblock Reconciling modern machine-learning practice and the classical bias-variance trade-off.
\newblock \emph{Proceedings of the National Academy of Sciences}, 116(32):15849--15854, 2019.

\bibitem[Nakkiran et~al.(2019)]{nakkiran2019deep}
P.~Nakkiran, G.~Kaplun, Y.~Bansal, T.~Yang, B.~Barak, and I.~Sutskever.
\newblock Deep double descent: Where bigger models and more data can hurt.
\newblock \emph{arXiv preprint arXiv:1912.02292}, 2019.

\bibitem[Advani and Saxe(2017)]{advani2017high}
M.~S. Advani and A.~M. Saxe.
\newblock High-dimensional dynamics of generalization error in neural networks.
\newblock \emph{arXiv preprint arXiv:1710.03667}, 2017.
\end{thebibliography}

\end{document}
